\documentclass[11pt]{article}
\usepackage[margin=1in]{geometry}
\usepackage{amsmath,amssymb,amsthm}
\newtheorem{theorem}{Theorem}
\newtheorem{proposition}{Proposition}
\newtheorem{corollary}{Corollary}
\newcommand{\Z}{\mathbb{Z}}
\newcommand{\vv}{v_2}
\DeclareMathOperator{\Imag}{Im}
\title{A 2-adic proof of the two-row $d=4$ fiber law for all $b\equiv 1\pmod 4$}
\author{Clio}
\date{2026-06-07}
\begin{document}
\maketitle

\noindent\textbf{Object.} For $\lambda=(2m-b,b)$, $0\le b\le m$, set $s=1+i$,
$P(u)=(1+su+u^2)^m$, and
\[
   I_b(m):=\Imag\,[u^b]\big((1-u)P(u)\big)\in\Z .
\]
A previous cycle proved that $I_b(m)$ is a polynomial in $m$ whose only integer roots are the
\emph{forced} roots $m\in\{0,\dots,R\}$, $R:=\lfloor (b-1)/2\rfloor$, and that, writing
$(m)_{R+1}:=\prod_{r=0}^{R}(m-r)$ and $I_b(m)=(m)_{R+1}Q_b(m)$, the two-row $d=4$ fiber law
$G_{(2m-b,b)}(i)=0\iff(2,2)$ is equivalent to
\[
   (\star)\qquad Q_b(m)\ne 0\quad\text{for every integer }m\ge b .
\]
Since $m\ge b>R$ gives $(m)_{R+1}\ne0$, $(\star)$ is the same as $I_b(m)\ne0$ for $m\ge b$.

\section{An exact binomial expansion}
Writing $1+su+u^2=(1+u^2)+s\,u$ and expanding,
$P(u)=\sum_{j\ge0}\binom{m}{j}s^ju^j(1+u^2)^{m-j}$, so
\[
   I_b(m)=\sum_{j=1}^{b}\tau_j(m),\qquad
   \tau_j(m)=\binom{m}{j}\Imag\!\big((1+i)^j\big)\big(T(b-j)-T(b-1-j)\big),
   \tag{E}
\]
where $T(l)=[u^l](1+u^2)^{m-j}=\binom{m-j}{l/2}$ for $l$ even and $0$ for $l$ odd. Exactly one
of $b-j,b-1-j$ is even, so $T(b-j)-T(b-1-j)=\pm\binom{m-j}{\beta_j}$, $\beta_j:=\lfloor(b-j)/2\rfloor$.
Moreover $\Imag((1+i)^j)=2^{j/2}\sin(j\pi/4)\in\Z$, with $\vv(\Imag((1+i)^j))=\lfloor j/2\rfloor$
for $j\not\equiv0\ (4)$ and $\Imag((1+i)^j)=0$ for $j\equiv0\ (4)$. In particular
$\Imag((1+i)^j)$ is \emph{odd iff $j=1$}.

\begin{corollary}[mod $2$ law]
For every $b$, $\ I_b(m)\equiv (R+1)\binom{m}{R+1}\pmod 2$.
\end{corollary}
\begin{proof}
Reduce (E) mod $2$; only $j=1$ survives, and
$|\tau_1(m)|=m\binom{m-1}{R}=(m)_{R+1}/R!=(R+1)\binom{m}{R+1}$.
\end{proof}

\section{The law for $b\equiv1\pmod4$}
\begin{theorem}
Let $b\equiv1\pmod4$, $b\ge5$, $R=(b-1)/2$. Then for every integer $m$,
\[
   \vv\big(I_b(m)\big)=\vv\big((m)_{R+1}\big)-\vv(R!).
   \tag{V}
\]
Hence $\vv(Q_b(m))=-\vv(R!)$ is constant, $Q_b(m)\ne0$ for all integer $m$, and the two-row
$d=4$ law holds for all $b\equiv1\pmod4$.
\end{theorem}
\begin{proof}
As $b$ is odd with $b-1=2R$ even and $b-2$ odd, the $j=1$ term of (E) is exactly
$\tau_1(m)=m\binom{m-1}{R}=(m)_{R+1}/R!$, so $\vv(\tau_1(m))$ equals the right side of (V).
We show $\vv(\tau_j(m))\ge\vv(\tau_1(m))+1$ for all $2\le j\le b$ (those $j\equiv0\ (4)$ give
$\tau_j\equiv0$ and may be ignored).

Fix such a $j$ and put $h:=\lfloor j/2\rfloor$, $d_j:=\lfloor(j-1)/2\rfloor$, $s_j:=j+\beta_j$.
For $b$ odd one has $s_j=\lfloor(b+j)/2\rfloor$, $s_j-j=R-h$, and $d_j=s_j-(R+1)$. From
$\binom{m}{j}\binom{m-j}{\beta_j}=\binom{m}{s_j}\binom{s_j}{j}$,
\[
   \vv(\tau_j(m))=\vv\!\binom{m}{s_j}+\vv\!\binom{s_j}{j}+\Big\lfloor\tfrac j2\Big\rfloor .
\]
Telescoping $\binom{m}{s_j}/\binom{m}{R+1}=\prod_{t=1}^{d_j}(m-R-t)/(R+1+t)$ and using that
$R+1=(b+1)/2$ is \emph{odd} (so $\vv(\tau_1)=\vv\binom{m}{R+1}$) gives
\[
   \vv(\tau_j(m))-\vv(\tau_1(m))=D_j+\sum_{t=1}^{d_j}\vv(m-R-t),\qquad
   D_j:=\vv\!\binom{s_j}{j}+\Big\lfloor\tfrac j2\Big\rfloor-\sum_{t=1}^{d_j}\vv(R+1+t).
   \tag{$\Delta$}
\]
With $\vv((2h)!)=\vv((2h+1)!)=h+\vv(h!)$ and
$\sum_{t=1}^{d_j}\vv(R+1+t)=\vv((R+1)!)-\vv((s_j-j)!)$ (as $s_j-j=R-h$), a direct simplification
collapses $D_j$ to a product of $h+1$ consecutive integers divided by $h!$:
\[
   D_j=\vv\big((h+1)!\,\tbinom{R+1}{h+1}\big)-\vv(h!)=\vv(h+1)+\vv\!\binom{R+1}{h+1}\ \ge\ 0.
   \tag{D}
\]
Now combine $(\Delta)$ and (D):
\begin{itemize}
\item If $D_j\ge1$, then $\vv(\tau_j)-\vv(\tau_1)\ge1$.
\item If $D_j=0$, then by (D) $h+1$ is odd ($h$ even) and $\binom{R+1}{h+1}$ is odd. Since
$j\equiv0\ (4)$ is excluded, $h$ even forces $j$ odd, i.e. $j=2h+1\equiv1\ (4)$ with $h\ge2$
(the case $h=0$ is $j=1$). Then $d_j=h\ge2$, and among the $\ge2$ consecutive integers
$m-R-1,\dots,m-R-d_j$ at least one is even, so $\sum_{t}\vv(m-R-t)\ge1$ and again
$\vv(\tau_j)-\vv(\tau_1)\ge1$.
\end{itemize}
Thus the tail $\sum_{j\ge2}\tau_j$ has $\vv\ge\vv(\tau_1)+1$, whence
$\vv(I_b(m))=\vv(\tau_1(m))$, which is (V).
\end{proof}

\section{The remaining classes}
For $b\equiv0\pmod4$, $R+1=b/2$ is even, so $D_j$ acquires a correction $-\vv(R+1)$ and the
$j=1$ term only \emph{ties} (never loses) in valuation; the identity (V) still holds
numerically (verified $b\le40$, random $m\le3\cdot10^4$), the sole gap being a mod-$4$
refinement of the Corollary showing the valuation-minimal terms sum to a $2$-adic unit. For
$b\equiv2,3\pmod4$ the monic integer part $q_b$ of $Q_b$ satisfies
$q_b(m)\equiv m\,(m^2+m+1)^{\lfloor(b-1)/4\rfloor}\pmod 2$, i.e.\ it has a genuine root
$m\equiv0\pmod2$; the $2$-adic method is structurally inapplicable and an odd-prime
Newton-polygon / irreducibility input is required. (Equivalently, for $b\equiv0,1\pmod4$,
$q_b\equiv(m^2+m+1)^{\lfloor(b-1)/4\rfloor}\pmod2$, which has no root mod $2$.)
\end{document}
